% SPDX-License-Identifier: CC-BY-SA-4.0
% Author: Matthieu Perrin
% Part: <Nom de la partie>
% Section: <Nom de la section>
% Sub-section: <Nom de la sous-section>  % (facultatif, laisser vide si non utilisé)
% Frame: <Titre de la slide>

\begingroup

\begin{frame}{Système d'exploitation temps-réel}

  \on[text, top=-3mm]{
    \Probleme{RT-Scheduling}{
      Un ensemble fini de tâches $J \subseteq \mathbb{N}^3$
      \begin{itemize}
      \item Chaque tâche est un triplet $j = \langle d_j, s_j, e_j \rangle \in J$ :
        \begin{description}[$d$ :]
        \item[$d_j$ :] durée de traitement
        \item[$s_j$ :] date de disponibilité
        \item[$e_j$ :] date d'échéance
        \end{description}
      \end{itemize}
    }{
      Existe-t-il un ordonnancement $o : J \rightarrow \mathcal{P}(\mathbb{N})$ tel que :
      \begin{itemize}
      \item $\forall j \in J,\exists t\in\mathbb{N},~ \structure{o(j) = [t, t + d_j - 1]}$
      \item $\forall j \in J,~\quad\quad\quad         \structure{o(j) \subseteq [s_j, e_j]}$
      \item $\forall j_1 \neq j_2 \in J,~\quad           \structure{o(j_1) \cap o(j_2) = \emptyset}$
      \end{itemize}
    }
  }

  \onBlock[y=-13mm, anchor=north]{Théorème (admis)}{
    \begin{itemize}
      \item \textsc{RT-Scheduling} est \textsc{np}-complet
      \item Version \structure{préemptive} dans \textsc{p}.
        \begin{itemize}
        \item Première condition remplacée par $|o(j)| = d_j$
        \end{itemize}
    \end{itemize}
  }

  \on[y=-25mm, x=30mm]{
    \begin{tikzpicture}[x=5mm, y=4mm]\tiny
      
      \foreach \x in {1,2,3}{
        \draw[densely dotted, black!50] (0.5,\x) -- (8.5,\x);
      }
      \foreach \x in {1,2,...,8}{
        \draw[densely dotted, black!50] (\x,0.5) -- (\x,3.5);
        \node[black!50] at (\x,.2) {\x};
      }

      \draw[alert, thick, -|] (0.5,3) -- (2,3);
      \draw[alert, thick, -|] (0.5,2) -- (1,2);
      \draw[alert, thick, -|] (0.5,1) -- (3,1);
      
      \draw[alert, thick, -|] (8.5,3) -- (6,3);
      \draw[alert, thick, -|] (8.5,2) -- (8,2);
      \draw[alert, thick, -|] (8.5,1) -- (7,1);

      \draw[example, fill=example!30, rounded corners=.5mm] (2.05,2.8) rectangle (3.95,3.2);
      \draw[example, fill=example!30, rounded corners=.5mm] (5.05,1.8) rectangle (7.95,2.2);
      \draw[example, fill=example!30, rounded corners=.5mm] (4.05,0.8) rectangle (4.95,1.2);

      \node[example, right] at (8.5,3) {$\langle 2, 2, 6 \rangle$};
      \node[example, right] at (8.5,2) {$\langle 3, 1, 8 \rangle$};
      \node[example, right] at (8.5,1) {$\langle 1, 3, 7 \rangle$};
      
    \end{tikzpicture}
  }
  
  \footnoteref{Garey, Johnson, Sethi. \emph{The Complexity of Flowshop and Jobshop Scheduling}. MOR (1976)}
  
\end{frame}

\endgroup
